\subsection{An integral functorial closure of an odd-to-even disk}

\begin{theorem}[D-029: \status{Proved}; the regraded divided-power bar closure]
\label{sa08:D029}
Let \(M\) be a free module over a commutative ring \(k\),
and let \(r\ge1\). There is a functorial augmented strictly
commutative PD cochain algebra \(U_r(M)\), with a second
grading called arity, and an injective morphism
\[
 \Gamma_k(M_y)\otimes_k\Lambda_k(M_x)
       \longrightarrow U_r(M),
 \qquad |x|=2r-1,\quad |y|=2r,\quad dx=y.
 \tag{D-29a}
\]
Both \(x\) and \(y\) have arity one. The augmentation
\(U_r(M)\to k\) is a cochain homotopy equivalence, by
an explicit natural integral contraction. Its arity-\(m\)
part is supported in form degrees
\[
                  (2r-1)m\le q\le2rm.
 \tag{D-29b}
\]
Every term is free over \(k\), and the construction and
contraction commute with extension of the coefficient ring.

If \(M\) is also a degreewise free simplicial module
whose normalized complex is acyclic, every positive-arity
form row of \(U_r(M)\) is flat and horizontally acyclic.
The same holds after tensoring such a row with any
simplicial module. Consequently tensoring an eventual
theory with this closure on an acyclic simplicial parameter
is an admissible absolute extension. This does not assert
a relative cell-attachment or arbitrary-target lifting theorem.
\end{theorem}

\begin{proof}
Put \(B=\Gamma_k(M)\), with augmentation ideal
\(\overline B\). As a module, take the one-sided normalized
bar construction
\[
 U_r(M)=\bigoplus_{n\ge0} B\otimes_k\overline B^{\otimes n}.
 \tag{D-29c}
\]
Write its elements as \(b_0[b_1|\cdots|b_n]\).
All the \(b_i\) retain their divided-power arities.
An element of total arity \(m\) and bar length \(n\)
has form degree \(q=2rm-n\). Since every bar entry
has positive arity, \(0\le n\le m\), proving (D-29b).

The differential is
\[
\begin{aligned}
d\bigl(b_0[b_1|\cdots|b_n]\bigr)
={}&b_0b_1[b_2|\cdots|b_n]\\
 &+\sum_{i=1}^{n-1}(-1)^i
 b_0[b_1|\cdots|b_i b_{i+1}|\cdots|b_n].
\end{aligned}
 \tag{D-29d}
\]
It raises form degree by one. Multiply coefficients
in \(B\), and use signed shuffles of bar entries for
the product. Each bar entry has odd regraded degree
\(2ra_i-1\); hence the signs are precisely the ordinary
permutation signs. The usual pairing of two consecutive
merges proves \(d^2=0\), and the shuffle identities prove
associativity, strict graded commutativity, and the
Leibniz rule. Equivalently, this is normalization of
the simplicial commutative bar algebra, followed by the
displayed regrading. The standard normalized bar
PD structure is discussed in
\cite[Theorem~7.5, pp.~250--251; \S8, pp.~252--253]{D_Richter2009}.

For completeness, its integrality here can be seen directly.
Work first over \(\Z\) with a basis of \(M\).
For an even bar-length basis word of positive length,
the signed shuffles in its \(a\)-th power occur in free
orbits under permutation of its \(a\) identical blocks.
The block permutations have positive sign. Thus division
of that power by \(a!\) is integral. For bar length zero,
the divided powers are those already present in \(B\).
Multiplying a positive-length word by a coefficient in
\(B\) gives the usual scalar rule for its divided powers.
The divided-power sum formula then defines the operations
on all even elements of positive arity. Odd squares
vanish by the sign-reversing interchange of the two
blocks. All PD identities and
\(d\gamma_a(z)=\gamma_{a-1}(z)\,dz\) for even \(z\)
hold after rationalization; the integral underlying
modules are free, so these identities hold integrally.
They consequently hold after arbitrary base change.
The same integral argument proves compatibility with
linear maps of the parameter module: first use a polynomial
ring over \(\Z\) with indeterminate matrix entries,
then specialize those entries to the chosen coefficients. Infinite free
modules reduce to finite subsets of their bases in
each expression.

Let \(\epsilon:B\to k\) be the augmentation and put
\(\overline b_0=b_0-\epsilon(b_0)\). Define
\[
 h\bigl(b_0[b_1|\cdots|b_n]\bigr)
       =1[\overline b_0|b_1|\cdots|b_n].
 \tag{D-29e}
\]
The right side is zero when \(b_0\) is a scalar.
The first merge in \(dh\) recovers the original
positive-arity element; its next merge cancels the
first term of \(hd\), and all remaining terms cancel
in consecutive pairs. Thus
\[
                 dh+hd=\mathrm{id}-\eta\epsilon,
                 \qquad h^2=0.
 \tag{D-29f}
\]
This is natural in \(M\), preserves arity, and lowers
form degree by one. It proves the augmented cochain
homotopy equivalence over every \(k\).

Send \(y(a)\) to the coefficient \(a\in B\), and
\(x(a)\) to \(1[a]\). The shuffle product sends an
exterior monomial to its alternating sum of words.
For a free basis of \(M\), each such sum contains a
distinguished word with coefficient \(1\), and distinct
exterior monomials have disjoint word supports. Therefore
the map (D-29a) is injective, with a degreewise split
underlying module inclusion. Formula (D-29d) gives
\(dx(a)=y(a)\); the construction preserves divided powers.

Finally, each fixed-arity term of (D-29c) is a finite
direct sum of tensor products
\[
 \Gamma^{a_0}(M)\otimes\Gamma^{a_1}(M)\otimes\cdots
       \otimes\Gamma^{a_n}(M),
 \qquad a_0\ge0,\quad a_i>0\ (i>0),\quad\sum a_i=m.
\]
For \(m>0\), at least one factor has positive arity.
Statement~\ref{sa08:C05} makes that factor flat and
universally horizontally acyclic, while all other
factors are flat. This proves the horizontal assertion.
The natural contraction (D-29e) commutes with the
outer simplicial maps, so tensoring with an existing
theory preserves its vertical resolution as well.
In a fixed form degree only finitely many arities occur
by (D-29b); the new positive-arity summands are
horizontally acyclic and the arity-zero summand is
the original theory. Assigning external weight \(mw\)
to arity \(m\), with \(w\ge2r\), also respects the
required inequality between form degree and weight.
\end{proof}

\paragraph{The first repair is already present integrally.}
For one parameter and \(r=1\), write
\[
 x=1[y],\quad z=1[y^{[2]}],\quad v=1[y|y].
\]
Then
\[
 dx=y,\qquad dz=y^{[2]},\qquad dv=xy-2z.
 \tag{D-29g}
\]
The new parameters \(z,v\) have form degrees three
and two, respectively. There is no new form-degree-one
parameter in arity two. The bar differential and the
contraction include every higher repair integrally,
without a choice of a prime.

The program \texttt{code/D\_bar\_closure\_check.py}
records all rank-one differential matrices through
arity eight, verifies \(d^2=0\) and (D-29f) on every
basis vector, and checks low-arity shuffle signs,
the Leibniz rule, odd squares, and divided squares.
Its \(2144\) checks and complete matrices are saved
in \texttt{logs/D\_bar\_closure\_check.json}.
Existence of the closure object does not itself give
a compatible map from it to a prescribed target;
Statement~\ref{sa08:D030} treats only its first nontrivial arity.
